\section{Matching and relation corners}
\label{sec:matching}

We first establish the algebraic properties used in the rank argument.
All monoid products in this paper are in path order: in a product $ab$,
the factor $a$ is placed to the left of the factor $b$.
A homomorphism is called \emph{unital} when it preserves the designated
identities.

\subsection{Matching diagrams, rank, and support}

For an integer $m\geq 0$, let $[m]=\{1,\ldots,m\}$, with
$[0]=\varnothing$.  A \emph{matching diagram of degree $m$} is a perfect
matching on two disjoint sets of ports, a left set and a right set, each
indexed by $[m]$.  Only the paired ports matter; crossings in a drawing
have no significance.  Let $\D{m}$ be the set of these diagrams.

To multiply $a,b\in\D{m}$, identify the right port indexed $i$ of $a$
with the left port indexed $i$ of $b$, for every $i\in[m]$.  In the
resulting finite multigraph, each outside port has degree one and each
identified port has degree two.  Thus every component meeting an outside
port is a path with two outside endpoints.  Pair these endpoints and
discard the components without outside ports.  The result is $ab$.
Parallel edges, when they occur at the interface, are counted separately.

This multiplication is associative.  Indeed, in any string of glued
diagrams, a substring can be replaced by the matching of its boundary
endpoints: each path through the substring is replaced by one edge with
the same endpoints.  This preserves the connections between the outer
ports of the whole string.  Components internal to the substring have
no connection to those ports.  Contracting the first or the last two
factors in a string of three therefore gives the same outside matching.
The diagram $\id_m$ pairing each left port $i$ with right port $i$ is
an identity.  In particular, $\D{0}=\{\id_0\}$ is the one-element
monoid.  The monoids $\D{m}$ are the Brauer diagram
monoids, whose matchings underlie Brauer's
algebras~\cite{Brauer1937}. Their semigroup structure and transversal
rank have been studied in particular by Auinger and by East and
Gray~\cite{Auinger2012,EastGray2017}.

A pair joining opposite sides is a \emph{transversal}; a pair on one
side is a \emph{cap}.  Write $\rank(a)$ for the number of transversals
of $a$, and define
\[
  \supp(a)=\{i\in[m]:
      \text{$a$ does not pair left $i$ with right $i$}\}.
\]
We say that $a$ is \emph{supported in} $J\subseteq[m]$ if
$\supp(a)\subseteq J$.

\begin{lemma}[Rank and support]
\label{lem:rank-basic}
For $a,b\in\D{m}$,
\[
  \rank(ab)\leq \min\{\rank(a),\rank(b)\},
  \qquad
  m-\rank(a)\leq |\supp(a)|.
\]
The units of $\D{m}$ are exactly its rank-$m$ diagrams, which are the
permutation diagrams.  The only rank-$m$ idempotent is $\id_m$.
Conjugation by a permutation diagram preserves rank.
For every $J\subseteq[m]$, the diagrams supported in $J$ form a
submonoid of $\D{m}$.
\end{lemma}

\begin{proof}
Each transversal path in the product uses a transversal of $a$ and a
transversal of $b$.  Different product paths are disjoint, so they use
distinct transversals of either factor.  This proves the rank inequality.
Every index outside $\supp(a)$ supplies a fixed transversal, proving
the support inequality.

A rank-$m$ diagram is a bijection between its two sides and has the
inverse permutation diagram as a two-sided inverse.  Conversely, the
rank inequality forces every unit to have rank $m$.  An idempotent in
a group is its identity, proving the assertion about idempotents.
If $u$ is a permutation diagram, the rank inequality gives
$\rank(uau^{-1})\leq\rank(a)$; applying it to
$a=u^{-1}(uau^{-1})u$ gives the reverse inequality.

If both $a$ and $b$ fix the transversal at an index $i$, their product
also fixes it.  Hence any product of diagrams supported in $J$ is
supported in $J$, and $\id_m$ is supported in every such set.
\end{proof}

\subsection{Reduction at an idempotent}

If $e^2=e$ in a monoid and $S$ is a subset closed under multiplication
and containing $e$, its \emph{corner}
\[
  eSe=\{ese:s\in S\}
\]
is a monoid with identity $e$.  In particular, every $z\in eSe$
satisfies $ez=ze=z$. This is the standard local-monoid construction
of finite-semigroup theory~\cite[Section~2.7]{Pin1997}.
For $m\ge3$, Auinger observes that the corner of the standard
rank-$(m-2)$ projection in $\D{m}$ is isomorphic to
$\D{m-2}$~\cite[Section~3]{Auinger2012}.
The following reduction treats an arbitrary idempotent and realizes its
matching corner on fewer ports. Its third part controls one common support
set for an entire family of sandwiches.

\begin{lemma}[Matching corners]
\label{lem:corners}
Let $e\in\D{m}$ be idempotent and let $r=\rank(e)$.
\begin{enumerate}[label=(\roman*)]
  \item There is a unital injective homomorphism
  \[
    R_e:e\D{m}e\longrightarrow\D{r}
  \]
  satisfying $\rank(R_e(z))=\rank(z)$ for all $z\in e\D{m}e$.
  \item The idempotent satisfies
  \[
    |\supp(e)|\leq 2(m-r).
  \]
  \item For every $J\subseteq[m]$ there is a set $J'\subseteq[r]$,
  depending only on $e$, $J$, and the numbering chosen for $R_e$, such
  that $|J'|\leq |J|$ and
  \[
    \supp(R_e(eae))\subseteq J'
    \quad\text{for every $a\in\D{m}$ supported in $J$}.
  \]
\end{enumerate}
\end{lemma}

\begin{proof}
Let $X,Y\subseteq[m]$ be the indices of the left and right transversal
endpoints of $e$, respectively.  Its transversals define a bijection
$p:X\to Y$, with $|X|=|Y|=r$.

We describe the connector at the interface between two copies of $e$.
Keep two distinct tagged rows of vertices, $R_i$ and $L_i$ for
$i\in[m]$, and put an edge $R_iL_i$ between each matching pair of
indices.  On the $R$ row put the right caps of $e$, and on the $L$ row
put the left caps of $e$.  Thus the degree-one vertices are precisely
\[
  \{R_y:y\in Y\}\ \cup\ \{L_x:x\in X\},
\]
and all other vertices have degree two.  These tagged vertices remain
distinct even when $X$ and $Y$ overlap.  Pairing the endpoints of its
paths, and discarding its closed components, gives a matching on the
$2r$ displayed vertices.

This is exactly the connector encountered by the transversals in
$e\id_m e=e^2$.  Since $\rank(e^2)=r$, all $r$ left transversal
endpoints of the first copy must reach the right boundary of the second.
Consequently the connector cannot pair two $R$ endpoints or two $L$
endpoints: it pairs $R_y$ with $L_{\gamma(y)}$ for a bijection
$\gamma:Y\to X$.  The path from an outer left endpoint $x\in X$
ends at the right endpoint $p(\gamma(p(x)))$.  Equality $e^2=e$ gives
\[
  p(\gamma(p(x)))=p(x),
  \qquad\text{hence}\qquad \gamma(p(x))=x.
\]
Thus $\gamma=p^{-1}$.

Every diagram $z\in e\D{m}e$ contains the left caps of $e$ on its
outer left side and the right caps of $e$ on its outer right side.
Delete precisely these prescribed caps.  The remaining ports are the
left ports indexed by $X$ and the right ports indexed by $Y$, and they
retain a perfect matching.  In particular, any additional caps among
these retained ports remain in the reduced diagram.  Number $X$ by
$[r]$ in any fixed way, and give the right endpoint $p(x)$ the same
number as $x$.  This defines $R_e(z)$.

Figure~\ref{fig:corner-reduction} illustrates the connector and the
numbering for a rank-$2$ idempotent.
\begin{figure}[htbp]
\centering
\begin{tikzpicture}[
  >=Stealth,
  every node/.style={font=\small},
  first/.style={very thick,blue!65!black},
  second/.style={very thick,orange!70!black},
  port/.style={circle,fill=gray!70,inner sep=1.4pt},
  retained/.style={circle,draw,fill=white,inner sep=1.7pt}
]
\begin{scope}
\node at (0.9,2.05) {idempotent $e\in\D{4}$};
\foreach \i/\y in {1/1.2,2/0.4,3/-0.4,4/-1.2} {
  \coordinate (eL\i) at (0,\y);
  \coordinate (eR\i) at (1.8,\y);
  \node[anchor=south] at ([yshift=2pt]eL\i) {$\i$};
  \node[anchor=south] at ([yshift=2pt]eR\i) {$\i$};
}
\draw[first] (eL1)--(eR1);
\draw[second] (eL2)--(eR4);
\draw[thick,gray] (eL3) .. controls (-0.55,-0.4) and (-0.55,-1.2) .. (eL4);
\draw[thick,gray] (eR2) .. controls (2.35,0.4) and (2.35,-0.4) .. (eR3);
\foreach \p in {eL3,eL4,eR2,eR3} \node[port] at (\p) {};
\foreach \p in {eL1,eL2,eR1,eR4} \node[retained] at (\p) {};
\node[align=center] at (0.9,-1.85)
  {$X=\{1,2\}$,\quad $Y=\{1,4\}$\\$p(1)=1,\quad p(2)=4$};
\end{scope}

\draw[->,thick] (2.7,0)--node[above] {glue} (3.4,0);

\begin{scope}[shift={(4.3,0)}]
\node at (0.9,2.05) {connector in $e\id_4e$};
\foreach \i/\y in {1/1.2,2/0.4,3/-0.4,4/-1.2} {
  \coordinate (cR\i) at (0,\y);
  \coordinate (cL\i) at (1.8,\y);
  \node[anchor=south] at ([yshift=2pt]cR\i) {$R_{\i}$};
  \node[anchor=south] at ([yshift=2pt]cL\i) {$L_{\i}$};
}
\draw[first] (cR1)--(cL1);
\draw[second] (cR4)--(cL4);
\draw[second] (cL4) .. controls (2.35,-1.2) and (2.35,-0.4) .. (cL3);
\draw[second] (cL3)--(cR3);
\draw[second] (cR3) .. controls (-0.55,-0.4) and (-0.55,0.4) .. (cR2);
\draw[second] (cR2)--(cL2);
\foreach \p in {cR2,cR3,cL3,cL4} \node[port] at (\p) {};
\foreach \p in {cR1,cR4,cL1,cL2} \node[retained] at (\p) {};
\node at (0.9,-1.85) {$\gamma(1)=1,\quad\gamma(4)=2$};
\end{scope}

\draw[->,thick] (7.0,0)--node[above,align=center] {retain and\\renumber} (8.15,0);

\begin{scope}[shift={(9.1,0)}]
\node at (0.9,2.05) {$R_e(e)=\id_2$};
\coordinate (dL1) at (0,0.6);
\coordinate (dL2) at (0,-0.6);
\coordinate (dR1) at (1.8,0.6);
\coordinate (dR2) at (1.8,-0.6);
\draw[first] (dL1)--(dR1);
\draw[second] (dL2)--(dR2);
\foreach \p in {dL1,dL2,dR1,dR2} \node[retained] at (\p) {};
\node[anchor=south] at ([yshift=2pt]dL1) {$1$};
\node[anchor=south] at ([yshift=2pt]dL2) {$2$};
\node[anchor=south] at ([yshift=2pt]dR1) {$1$};
\node[anchor=south] at ([yshift=2pt]dR2) {$2$};
\node[align=center] at (0.9,-1.85)
  {$x$ and $p(x)$ share a number;\\right $4=p(2)$ becomes $2$};
\end{scope}
\end{tikzpicture}
\caption{A rank-$2$ idempotent of degree $4$ and its reduction. In the
left panel, the gray caps are prescribed and the colored transversals
are retained. The connector paths are $R_1$--$L_1$ and
$R_4$--$L_4$--$L_3$--$R_3$--$R_2$--$L_2$, giving
$\gamma=p^{-1}$. Numbering left $x$ and right $p(x)$ alike yields
$R_e(e)=\id_2$. For a general corner element $z$, only the prescribed
outer caps are deleted; every pair among retained ports, including
caps, remains.}
\label{fig:corner-reduction}
\end{figure}


Restoring the prescribed caps and the original indices reconstructs
$z$, so $R_e$ is injective.  Only nontransversal pairs were deleted,
so it preserves rank, and by construction $R_e(e)=\id_r$.
To check multiplication, let $z,w\in e\D{m}e$.  The fixed caps at
their middle interface are the right caps and left caps of $e$.
Their connector is therefore the one just computed: the retained right
port $p(x)$ of $z$ is connected to the retained left port $x$ of $w$.
In the reduced numbering these ports have the same index.
Contracting the connector, including when the reduced diagrams have
caps of their own, thus gives ordinary multiplication in $\D{r}$.
It follows that $R_e(zw)=R_e(z)R_e(w)$, proving (i).

For (ii), suppose that $i\in X\cap Y$.  The vertices $R_i,L_i$ in
the connector have no incident caps, so their edge $R_iL_i$ is an
entire connector path.  Hence $\gamma(i)=i$, and $\gamma=p^{-1}$
implies $p(i)=i$.  Thus every index in $X\cap Y$ is fixed by $e$,
and
\[
  \supp(e)\subseteq([m]\setminus X)\cup([m]\setminus Y).
\]
Both complements have size $m-r$, which proves the bound.

For (iii), consider the full glued graph $e\id_m e$, keeping the
middle identity diagram explicit.  Its transversal paths have the same
outside endpoint pairs as $e$, and so are indexed by the $r$ numbers
used in the reduction.  Mark a transversal path if it uses any middle
identity edge whose index belongs to $J$.  Let $J'$ be the numbers of
the marked paths.  A transversal path may use several middle identity
edges; the only counting fact needed is that different paths are
disjoint.  Each middle identity edge therefore marks at most one
transversal path, and $|J'|\leq|J|$.

Now replace the middle identity by any diagram $a$ supported in $J$.
Every middle identity edge at an index outside $J$ remains present.
An unmarked transversal path used only such edges, so that entire path
survives with the same outside endpoints.  Each of its internal vertices
still has its two original incident edges; a new path cannot attach to
it.  In particular, changes to closed components or to paths with both
endpoints on the same side cannot join an unmarked transversal: all
its middle ports are already paired by the unchanged identity edges.
Its corresponding reduced port is therefore still paired with the
equally indexed port on the other side.  Only numbers in $J'$ can belong
to $\supp(R_e(eae))$.  The construction of $J'$ did not depend on $a$,
as required.

When $r=0$, the prescribed caps exhaust all outside ports.  Then
$eae=e$ for every $a$, and the corner is identified with the singleton
$\D{0}$.  The reduction above is its unique map, and (iii) holds with
$J'=\varnothing$; (ii) also holds since $|\supp(e)|\leq m$.
\end{proof}

The reduction will also be applied to a submonoid $eSe$.  We write
$R_e(eSe)$ for its image in $\D{r}$.  It contains the ambient identity
$\id_r$.  When discussing support in this reduced monoid, the port
set is $[r]$; support in the original monoid still refers to $[m]$.

\subsection{Homomorphic images and permutation lifts}

Every element of a finite semigroup has an idempotent positive power.
Indeed, for an element $x$ there are integers $a\geq1$ and $d\geq1$
such that $x^{k+d}=x^k$ for all $k\geq a$.  Choose $N\geq a$
divisible by $d$.  Then $x^{2N}=x^N$.  If a multiplicative map sends
$x$ to an idempotent $q$, the same is true of every positive power of
$x$, including this idempotent power.

\begin{lemma}[Unit lifts from a minimum-rank corner]
\label{lem:unit-lifts}
Let $S\subseteq\D{m}$ be nonempty and closed under multiplication,
and let $\phi:S\to M$ be an onto multiplicative map to a monoid $M$.
Let $q\in M$ be idempotent.
There is an idempotent $e\in S$ with $\phi(e)=q$.  For every such
$e$, the reduced monoid $R_e(eSe)$ has a unital surjection
\[
  \phi_e:R_e(eSe)\longrightarrow qMq,
  \qquad \phi_e(R_e(z))=\phi(z).
\]
If $e$ has minimum rank among all idempotents in $S$ mapping to $q$,
then every element of $eSe$ mapping to a unit of $qMq$ reduces to a
permutation diagram.  In particular, every unit of $qMq$ has a
permutation lift in $R_e(eSe)$, and that lift is a unit of the reduced
monoid.
\end{lemma}

\begin{proof}
The sets $S$ and $M$ are finite because $\D{m}$ is finite and $\phi$
is onto.  Take any preimage of $q$ and then an idempotent positive power
to obtain $e$.  The image of $eSe$ under $\phi$ is exactly $qMq$:
for any $v\in M$, a preimage $s\in S$ gives
$\phi(ese)=qvq$.  Its identity $e$ maps to the identity $q$ of this
corner.  Lemma~\ref{lem:corners}(i) makes the displayed definition of
$\phi_e$ well-defined and transfers both multiplication and surjectivity.

Suppose now that $e$ has the stated minimum rank, put $r=\rank(e)$,
and let $y\in eSe$ map to a unit $u$ of $qMq$.  The group of units
of $qMq$ is finite, so $u^t=q$ for some positive integer $t$.
Take an idempotent positive power $z$ of $y^t$.  Then $z\in S$ is
idempotent and $\phi(z)=q$.  The rank inequalities and minimum choice
of $e$ give
\[
  r\leq\rank(z)\leq\rank(y)\leq r.
\]
The last inequality uses $y\in eSe$.  Hence $R_e(y)$ has rank $r$
in $\D{r}$ and is a permutation diagram.  Its inverse is a power of
the same finite-order permutation, so lies in $R_e(eSe)$.  Surjectivity
of $\phi_e$ supplies a lift for each unit $u$ and completes the proof.
\end{proof}

\subsection{Corners of the relation monoid}

For a finite set $H$ and a subset $F\subseteq H$, we regard a relation
on $F$ as a relation on $H$ containing no pairs outside $F\times F$.
With $i=\Id{F}$, these are precisely the relations $A$ satisfying
$A=iAi$. Under this identification their monoid is $i\R{H}i$ with
identity $i$. Permutations of a finite set are identified with their
graph relations.

\begin{lemma}[A relation corner]
\label{lem:relation-corner}
Let $P\in\R{H}$ be idempotent, and let $F\subseteq H$ satisfy
$\Id{F}P\Id{F}=\Id{F}$.  Put $i=\Id{F}$ and $Q=PiP$.
Then $Q$ is idempotent, $PQ=QP=Q$, and restriction is a unital
isomorphism
\[
  \theta:Q\R{H}Q\longrightarrow\R{F},
  \qquad \theta(Z)=iZi.
\]
Its inverse sends $A\in\R{F}$ to $PAP=QAQ$.
\end{lemma}

\begin{proof}
The identities $P^2=P$ and $iPi=i$ give
\[
  Q^2=PiP^2iP=P(iPi)P=Q,
  \qquad PQ=QP=Q.
\]
They also give $iQ=iP$, $Qi=Pi$, and $iQi=i$.
If $Z\in Q\R{H}Q$, then $Z=QZQ$ and $PZP=Z$.  Consequently
\[
  Z=PiPZPiP=PiZiP.                         \tag{$*$}
\]
Thus its restriction $iZi$ determines $Z$ uniquely.

Conversely, if $A\in\R{F}$, so that $A=iAi$, then
\[
  QAQ=QiAiQ=PiAiP=PAP,
  \qquad
  i(PAP)i=iPiAiPi=A.
\]
Hence every relation on $F$ is a restriction of an element of the
corner, and $A\mapsto PAP$ is the inverse bijection by~$(*)$.

For multiplication, write $Z=PAP$ and $W=PBP$, where
$A=iZi$ and $B=iWi$ are relations on $F$.  Since $A=Ai$ and
$B=iB$, we have
\[
  ZW=PAP^2BP=PAPBP
     =PA(iPi)BP=PABP.
\]
Restricting this equality gives $\theta(ZW)=AB=\theta(Z)\theta(W)$.
Finally, the identity $Q$ of the source corner restricts to $i$, the
identity of $\R{F}$.  This proves unitality and the claimed isomorphism.
\end{proof}
